% Generated draft: requires mathematical and source review.
\documentclass[11pt]{article}
\usepackage[margin=25mm]{geometry}
\usepackage{fontspec}
\setmainfont{Latin Modern Roman}
\usepackage{amsmath,amssymb,mathtools,unicode-math}
\setmathfont{Latin Modern Math}
\usepackage{enumitem,needspace,xurl,hyperref}
\setlength{\parindent}{0pt}
\setlength{\parskip}{5pt}
\setlength{\emergencystretch}{3em}
\newcommand{\N}{\mathbb{N}}
\newcommand{\Z}{\mathbb{Z}}
\newcommand{\Q}{\mathbb{Q}}
\newcommand{\R}{\mathbb{R}}
\newcommand{\C}{\mathbb{C}}
\newcommand{\F}{\mathbb{F}}
\newcommand{\A}{\mathbb{A}}
\newcommand{\PP}{\mathbb{P}}
\newcommand{\E}{\mathbb{E}}
\newcommand{\eps}{\varepsilon}
\newcommand{\dd}{\,\mathrm{d}}
\begin{document}
\section*{UnsolvedMath ID: 30002481}
\textbf{Source problem code:} catalogue entry
\subsection{Definitions and mathematical statement}
The supplied question concerns a provability implication over a theory denoted by \(Z_3\), called third-order arithmetic. Its intended axiom system is not specified. Third-order arithmetic ordinarily has objects representing natural numbers, sets of natural numbers, and collections of sets of natural numbers; this description alone does not fix its comprehension or other axioms.

For the usual interpretation of determinacy, write \(\N^{\N}\) for Baire space and \(\N^{<\N}\) for the finite sequences of natural numbers. Given a payoff set \(A\subseteq\N^{\N}\), players I and II alternately choose natural numbers, with I moving first, producing \(x\in\N^{\N}\). Player I wins exactly when \(x\in A\); otherwise player II wins. A strategy assigns a natural-number move to each finite position at which the relevant player moves. It is winning if every play consistent with it is won by that player. The set \(A\) is determined if one of the players has a winning strategy.

In the usual descriptive-set-theoretic interpretation, a \(\Pi^1_1\) payoff set has a definition of the form
\[
 x\in A\quad\Longleftrightarrow\quad
 \forall y\in\N^{\N}\;\theta(x,y,p),
\]
where \(\theta\) is arithmetical and \(p\) denotes any permitted real parameters. The assertion \(\Pi^1_1\text{-}\mathrm{Det}\) says that every payoff set in the specified \(\Pi^1_1\) class is determined. The source does not specify whether real parameters are permitted, nor whether this is formalized by a formula scheme or by quantification over third-order payoff sets.

The notation \(0^{\#}\) conventionally denotes zero sharp, a real associated with the theory of the constructible universe \(L\) with indiscernibles. The source does not give a formula in third-order arithmetic expressing its existence, or specify how the necessary constructibility and indiscernibility notions are represented there.

With these unresolved choices explicitly retained, the intended question is whether
\[
 Z_3\ \vdash\
 \bigl(\Pi^1_1\text{-}\mathrm{Det}\ \Longrightarrow\
 \exists r\subseteq\N\;(r\text{ codes }0^{\#})\bigr).
\]
This display is a template for the intended implication, not a fully specified sentence of a fixed formal theory. A precise provability question requires the axioms of \(Z_3\), the determinacy formulation and parameter convention, and the internal predicate for coding \(0^{\#}\) to be supplied.
\subsection{Short English statement}
The question asks whether third-order arithmetic can prove that determinacy of all \(\Pi^1_1\) games implies the existence of zero sharp. The supplied record leaves the base theory and the internal formulations insufficiently specified, so the displayed implication remains a statement template requiring human mathematical review.
\subsection{Sources}
\begin{itemize}
\item Local UnsolvedMath registry, \texttt{data/unsolvedmath/problems.json}, permanent ID \(30002481\), source problem code \texttt{catalogue entry}. The supplied record belongs to the Oberwolfach Reports Open Problems collection and lists proposed year 2014; no individual proposer is identified.
\item Supplied report citation: \emph{Set Theory} (2014), Oberwolfach Reports, pp. 91--144, DOI \texttt{10.4171/owr/2014/02}. \url{https://doi.org/10.4171/owr/2014/02}. The supplied citation gives a report-wide page range, without an exact location for this question.
\item Supplied literature source: Oberwolfach set theory report. \url{https://ems.press/content/serial-article-files/46488}.
\item Supplied literature source: Cheng--Schindler, higher-order arithmetic. \url{https://arxiv.org/abs/1503.04000}.
\item Literature assessment dated exactly \texttt{2026-08-21}: the supplied curation labels the problem open, states that the classical implication from \(\Pi^1_1\) determinacy to the existence of \(0^{\#}\) is a theorem over \(\mathrm{ZF}\), and describes the higher-order-arithmetic formulation as imprecise and the corresponding \(Z_3\) question as unresolved. This is supplied dated curation derived from the OWR record, not an independently verified literature assessment. Its time-sensitive status and cited claims require independent verification.
\item The supplied registry assigns the statement the label \texttt{corrected\_verified} and reports a manual semantic, typographic, and LaTeX review of its cleanup. These supplied editorial labels do not resolve the formalization choices identified above or establish mathematical correctness.
\end{itemize}
\end{document}
